Deep generative models based on differential equations, such as diffusion models and Poission flow generative models, have emerged as powerful tools for modeling high-dimensional data, from image synthesis~\cite{Song2020ScoreBasedGM, Ho2020DenoisingDP, Karras2022ElucidatingTD, Xu2022PoissonFG, Xu2023PFGMUT} to biological data~\cite{antipfgm, Watson2022BroadlyAA}. These models use iterative backward processes that gradually transform a simple distribution~(\eg Gaussian in diffusion models) into a complex data distribution by solving a differential equations. The associated vector fields~(or drifts) driving the evolution of the differential equations are predicted by neural networks. The resulting sample quality can be often improved by enhanced simulation techniques 
%of the differential equations, 
but at the cost of longer sampling times.

% Deep generative models involving differential equations, such as diffusion models and Poisson flow generative models, have emerged as powerful tools for modeling high-dimensional data, including image synthesis~\cite{Song2020ScoreBasedGM, Ho2020DenoisingDP, Karras2022ElucidatingTD, Xu2022PoissonFG, Xu2023PFGMUT} and biological data modeling~\cite{antipfgm, Watson2022BroadlyAA}. These models use iterative backward processes, transforming a simple distribution (e.g., Gaussian in diffusion models) into a complex data distribution via ODE or SDE, with vector fields (or drifts) predicted by neural networks. Sample quality generally improves with enhanced simulation techniques, but at the cost of longer sampling times.

% The simulation accuracy of the differential equationhighly correlated with the There is a trade-off between the sampling errors and the sampling time when simulating the differential equations, as increased sampling time often results in a smaller sampling errors~(better sample quality).

Prior samplers for simulating these backward processes can be categorized into two groups: ODE-samplers whose evolution beyond the initial randomization is deterministic, and SDE-samplers where the generation trajectories are stochastic. Several works~\cite{Song2020ScoreBasedGM, JolicoeurMartineau2021GottaGF, Karras2022ElucidatingTD} show that these samplers demonstrate their advantages in different regimes, as depicted in \Figref{fig:main-line}. ODE solvers~\cite{Song2020DenoisingDI, lu2022dpm, Karras2022ElucidatingTD} result in smaller discretization errors, allowing for decent sample quality even with larger step sizes (\ie fewer number of function evaluations~(NFE)). However, their generation quality plateaus rapidly. In contrast, SDE achieves better quality in the large NFE regime, albeit at the expense of increased sampling time. To better understand these differences, we theoretically analyze SDE performance: the stochasticity in SDE contracts accumulated error, which consists of both the discretization error along the trajectories as well as the approximation error of the learned neural network relative to the ground truth drift (\eg score function in diffusion model~\cite{Song2020ScoreBasedGM}). The approximation error dominates when NFE is large (small discretization steps), explaining the SDE advantage in this regime. Intuitively, the stochastic nature of SDE helps "forget" accumulated errors from previous time steps.

Inspired by these findings, we propose a novel sampling algorithm called \textit{Restart}, which combines the advantages of ODE and SDE. As illustrated in \Figref{fig:main-demo}, the Restart sampling algorithm involves $K$ repetitions of two subroutines in a pre-defined time interval: a \textit{Restart forward process} that adds a substantial amount of noise, akin to "restarting" the original backward process, and a \textit{Restart backward process} that 
%reverts the previous step by running 
runs the backward ODE. The Restart algorithm separates the stochasticity from the drifts, and the amount of added noise in the Restart forward process is significantly larger than the small single-step noise interleaving with drifts in previous SDEs such as \cite{Song2020ScoreBasedGM, Karras2022ElucidatingTD}, thus amplifying the contraction effect on accumulated errors. By repeating the forward-backward cycle $K$ times, the contraction effect introduced in each Restart iteration is further strengthened. The deterministic backward processes allow Restart to reduce discretization errors, thereby enabling step sizes comparable to ODE. To maximize the contraction effects in practice, we typically position the Restart interval towards the end of the simulation, where the accumulated error is larger. Additionally, we apply multiple Restart intervals to further reduce the initial errors in more challenging tasks.

Experimentally, Restart consistently surpasses previous ODE and SDE solvers in both quality and speed over a range of NFEs, datasets, and pre-trained models. Specifically, Restart accelerates the previous best-performing SDEs by $10\times$ fewer steps for the same FID score on CIFAR-10 using VP~\cite{Song2020ScoreBasedGM} ($2\times$ fewer steps on ImageNet $64\times 64$ with EDM~\cite{Karras2022ElucidatingTD}), and outperforms ODE solvers even in the small NFE regime. When integrated into previous state-of-the-art pre-trained models, Restart further improves performance, achieving FID scores of 1.88 on unconditional CIFAR-10 with PFGM++~\cite{Xu2023PFGMUT}, and 1.36 on class-conditional ImageNet $64\times 64$ with EDM. To the best of our knowledge, these are the best FID scores obtained on commonly used UNet architectures for diffusion models without additional training. We also apply Restart to the practical application of text-to-image Stable Diffusion model~\cite{Rombach2021HighResolutionIS} pre-trained on LAION $512\times 512$. Restart more effectively balances text-image alignment/visual quality~(measured by CLIP/Aesthetic scores) and diversity~(measured by FID score) with a varying classifier-free guidance strength, compared to previous samplers.

Our contributions can be summarized as follows: \textbf{(1)} We investigate ODE and SDE solvers and theoretically demonstrate the contraction effect of stochasticity via an upper bound on the Wasserstein distance between generated and data distributions~(Sec~\ref{sec:dilemma}); \textbf{(2)} We introduce the Restart sampling, which better harnesses the contraction effect of stochasticity while allowing for fast sampling. The sampler results in a smaller Wasserstein upper bound (Sec~\ref{sec:method}); \textbf{(3)} Our experiments are consistent with the theoretical bounds and highlight Restart's superior performance compared to previous samplers on standard benchmarks in terms of both quality and speed. Additionally, Restart improves the trade-off between key metrics on the Stable Diffusion model~(Sec~\ref{sec:exp}).


% \begin{enumerate}
%     \item intro
%     \item background: (1) trade-off between SDE and ODE; (2) why the absolute performance of SDE is better than ODE when using larger NFE
%     \item method description (restart) motivated by the success of SDE. 
%     \item method analysis, with toy experiments
%     \item experiments
% \end{enumerate}
\begin{figure*}
    \centering
    \subfigure[]{\includegraphics[width=0.55\textwidth]{img/main_demo.png}\label{fig:main-demo}}
    \subfigure[]{\includegraphics[width=0.40\textwidth]{img/main_line.png}\label{fig:main-line}}
    \vspace{-3pt}
    \caption{\textbf{(a)} Illustration of the implementation of drift and noise terms in ODE, SDE, and Restart. \textbf{(b)}
    Sample quality versus number of function evaluations~(NFE) for different approaches. ODE~(\textcolor{black!70!green}{Green}) provides fast speeds but attains only mediocre quality, even with a large NFE. SDE~(\textcolor{orange!50!yellow}{Yellow}) obtains good sample quality but necessitates substantial sampling time. In contrast to ODE and SDE, which have their own winning regions, Restart~(\textcolor{black!30!red}{Red}) achieves the best quality across all NFEs.}
    \label{fig:main}
    \vspace{-12pt}
\end{figure*}
